
All numbers reported here derive from the 33 S2AG-resolved papers from the huashen218 proxy corpus, executed against the live S2AG API with real data (no mock data detected; \texttt{mock\textbackslash \_data\textbackslash \_detection: false} in experiment log).

\subsubsection{5.1 RQ1: Pipeline Infrastructure Validated (23/23 Tests Pass)}

The complete S2AG bibliometric pipeline passes all 23 unit tests in 0.85 seconds with zero failures. The test suite covers corpus ingestion, S2AG resolution, classification (all three schemes), graph construction, gate evaluation, and figure generation.

Test breakdown:
\begin{itemize}
\item \texttt{TestClassifyPaper}: 11 tests --- all pass
\item \texttt{TestClassifyAll}: 1 test --- pass
\item \texttt{TestEvaluateGate}: 6 tests --- all pass
\item \texttt{TestExtractPaperIds}: 4 tests --- all pass
\item \texttt{test\textbackslash \_pipeline\textbackslash \_smoke}: 1 test --- pass
\end{itemize}

The gate failure reported in Section 5.3 is a data access failure, not a pipeline failure. Every algorithmic component is validated. The infrastructure contribution is complete and independent of the corpus scale limitation.

\subsubsection{5.2 RQ2: FoS-Primary Classification Achieves 100\% Coverage versus 12\% for Venue Strings}

\begin{table*}[htbp]
\centering
\resizebox{\dimexpr 2\columnwidth+\columnsep\relax}{!}{%
\begin{tabular}{lllllll}
\toprule
Scheme & Type & Papers Classified & Coverage & ML\_NLP & HCI & Unclassified \\
\midrule
Scheme 1 & Venue string (broad) & 4 / 33 & 12.1\% & 3 & 1 & 29 \\
Scheme 2 & Venue string (narrow) & 4 / 33 & 12.1\% & 3 & 1 & 29 \\
\textbf{Scheme 3} & \textbf{FoS-primary} & \textbf{33 / 33} & \textbf{100.0\%} & \textbf{28} & \textbf{5} & \textbf{0} \\
\bottomrule
\end{tabular}
}
\end{table*}

Schemes 1 and 2 produce identical coverage because all classifiable papers have abbreviated venue names in the target sets; the 29 unclassified papers all have full proceedings names in S2AG that do not match any substring.

The root cause is that S2AG's \texttt{venue} field stores the full proceedings name as submitted by the publisher. Substring matching against ''NeurIPS'' fails silently on ''Proceedings of the 37th Annual Conference on Neural Information Processing Systems.'' The \texttt{fieldsOfStudy} field is format-independent and consistently populated by S2AG's knowledge graph. The 12\% coverage rate is a format mismatch artifact, not a corpus composition artifact.

This finding generalizes: any bibliometric study of an interdisciplinary corpus using S2AG data and venue-string classification will silently drop the majority of papers unless FoS-primary classification or full-proceedings-name normalization is used.

\subsubsection{5.3 RQ3: Directional Citation Signal --- Ratio $\approx$ 0.107, Gate Thresholds Not Met}

\textbf{Coverage and Gate Evaluation:}

\begin{itemize}
\item S2AG resolution rate: 33 / 49 = \textbf{67.3\%} (gate threshold: 70\%) --- gate FAIL by 2.7 percentage points
\item Cross-group edge counts per scheme:
\end{itemize}

\begin{table*}[htbp]
\centering
\resizebox{\dimexpr 2\columnwidth+\columnsep\relax}{!}{%
\begin{tabular}{lllll}
\toprule
Scheme & $ML\_NLP\rightarrow HCI$ & $HCI\rightarrow ML\_NLP$ & Total Cross-Group & Gate $\geq$ 30? \\
\midrule
Scheme 1 & 0 & 0 & 0 & FAIL \\
Scheme 2 & 0 & 0 & 0 & FAIL \\
Scheme 3 & 5 & 4 & \textbf{9} & FAIL \\
\bottomrule
\end{tabular}
}
\end{table*}

Maximum cross-group edges across all schemes: 9 (gate threshold: 30) --- gate FAIL by 21 edges.

Both pre-registered gate thresholds are not satisfied. Statistical tests (P1: chi-squared on $2\times 2$ contingency table, P2: proportion z-test, P3: betweenness centrality) are not executed. All three predictions are INCONCLUSIVE --- not REFUTED.

\textbf{Directed Edge Matrix (Scheme 3):}

\begin{table}[htbp]
\centering
\resizebox{\columnwidth}{!}{%
\begin{tabular}{llll}
\toprule
 & $\rightarrow$ ML\_NLP & $\rightarrow$ HCI & Outgoing Total \\
\midrule
\textbf{ML\_NLP $\rightarrow$} & 42 & 5 & 47 \\
\textbf{HCI $\rightarrow$} & 4 & 0 & 4 \\
\bottomrule
\end{tabular}
}
\end{table}

Citation proportions: $ML\_NLP\rightarrow HCI$ = 5/47 = \textbf{10.6\%}; $HCI\rightarrow ML\_NLP$ = 4/4 = \textbf{100.0\%}; ratio r = (5/47) / (4/4) = \textbf{0.107}.

Interpretation: The directional pattern is consistent with the asymmetric epistemic dependency hypothesis. ML/NLP alignment papers allocate 42 of 47 within-corpus outgoing citations (89.4\%) to other ML/NLP papers, with only 5 citations (10.6\%) crossing to HCI. The 5 HCI-classified papers resolved under Scheme 3 allocated all 4 outgoing within-corpus citations to ML/NLP work; no HCI-to-HCI within-corpus citation edges are present ($HCI\rightarrow HCI$ = 0).

An important caveat: at N = 9 cross-group edges total, this ratio is not statistically confirmable. No falsifier was triggered, but the observed pattern does not constitute statistical confirmation of H-CitAsym. These figures are presented as a preliminary directional estimate only.

\textbf{Unexpected Finding: $HCI\rightarrow HCI$ = 0}

The complete absence of HCI-to-HCI within-corpus citations was not anticipated even under H-CitAsym. At N = 5 HCI papers in the resolved set, this could be a sample artifact. The pattern suggests that HCI alignment papers in this proxy treat ML/NLP alignment work as their primary reference frame within the corpus, rather than prior HCI alignment work. This finding requires verification at larger corpus scale.

\subsubsection{5.4 Corpus Proxy Analysis}

The corpus extraction funnel: approximately 130 linked papers in the GitHub reading list $\rightarrow$ 49 extractable identifiers (37.7\% extraction rate) $\rightarrow$ 33 S2AG-resolved papers (67.3\% resolution rate of extracted IDs; approximately 8.3\% of the target 400-paper corpus).

All 16 unresolved papers are ACM Digital Library publications without arXiv preprints. This dropout is non-systematic with respect to venue group: both ML/NLP and HCI papers appear in the unresolved set, introducing no directional bias into the citation direction measurement.

\textbf{Summary of Results:}

\begin{table}[htbp]
\centering
\resizebox{\columnwidth}{!}{%
\begin{tabular}{lll}
\toprule
Research Question & Result & Confidence \\
\midrule
RQ1: Pipeline validated? & 23/23 tests pass & HIGH \\
RQ2: FoS-primary coverage? & 100\% vs. 12\% (Schemes 1/2) & HIGH \\
RQ3: Gate thresholds met? & No (coverage = 67.3\%, edges = 9) & HIGH \\
RQ3: Directional ratio? & 0.107 (consistent with H-CitAsym) & MEDIUM (N = 9) \\
P1/P2/P3 statistical tests? & INCONCLUSIVE --- gate not met & --- \\
\bottomrule
\end{tabular}
}
\end{table}

